% Java, C#, Smalltalk 

\section{Metrics Evaluation}

\subsection{Metrics Evaluation for C\#}

C\# is a statically typed, object-oriented programming language standardised as ECMA-334 and commonly used in conjunction with the Common Language Runtime (CLR) \cite{ecma334}. The type system and object model include nominal typing, interface-based abstraction, and managed runtime semantics.

In this section, we evaluate the language against the defined metrics, grounding each assessment in the formal specification and relevant academic literature.

\subsubsection{Type System and Subtyping Model}

\textbf{1) Nominal-Structural Typing Score $S_{nom}(L)$}

C\# employs a strictly nominal type system in which subtypes are explicitly declared using type names, inheritance relationships, and the implementation of the interface \cite{ecma334}. Structural typing is not supported at the language level.

$I_{struct}(\text{C\#}) = 0, I_{nom}(\text{C\#}) = 1 \therefore S_{nom}(\text{C\#}) = \frac{0 + 1}{2} = 0.5 $

\textbf{2) Variance Formalization Score $S_{var}(L)$}

C\# supports generics with explicit variance annotations (\texttt{in}, \texttt{out}) for interfaces and delegates, introduced in C\# 4.0. These are formally constrained to preserve type safety.  However, it only applies to interfaces and delegates but does not apply to classes.

$V_{exp}(\text{C\#}) = 1, V_{sound}(\text{C\#}) = 1 \therefore S_{var}(\text{C\#}) = 1 $

\textbf{3) Inheritance-Subtyping Separation Score $S_{rel}(L)$}

First, C\# explicitly supports subtyping independently of class inheritance through interface implementation: a class or struct may implement one or more interfaces without inheriting implementation, thereby introducing subtyping relations decoupled from class inheritance.

$I_{dec}(\text{C\#})=1$

Second, subtyping rules in C\# are not strictly derived from inheritance semantics. The type system defines subtyping via multiple mechanisms, including interface implementation, variance in generics, and implicit reference conversions, indicating independence from purely inheritance-based relations.

$I_{ind}(\text{C\#})=1$

Third, although the C\# specification distinguishes between inheritance (class derivation) and interface-based subtyping, this distinction is not fully formalized as separate mathematical relations but rather described operationally within the type system rules.

$I_{form}(\text{C\#})=0.5$

\[
    S_{rel}(\text{C\#}) = \frac{1 + 1 + 0.5}{3} \approx 0.83
\]

\textbf{Final metric value:}
\begin{flalign*}
    & \qquad TSSI(\text{C\#}) = 0.25 \cdot S_{nom}(\text{C\#}) + 0.3 \cdot S_{var}(\text{C\#}) + 0.45 \cdot S_{rel}(\text{C\#}) &&\\
    & \qquad = 0.25 \cdot 0.5 + 0.3 \cdot 1 + 0.45 \cdot \frac{5}{6} &&\\
    & \qquad \approx 0.8 &&
\end{flalign*}


\subsubsection{Abstraction and Encapsulation Mechanisms}

\textbf{1) Visibility Granularity Index $V(L)$}

C\# defines seven access modifiers: \texttt{private}, \texttt{protected}, \texttt{internal}, \texttt{protected internal}, \texttt{private protected}, \texttt{public}, \texttt{file} \cite{ecma334}.

\[
    V(\text{C\#}) = \frac{7}{7}
\]


\textbf{2) Module Isolation Strength $M(L)$}

Namespaces do not enforce strict encapsulation (no export lists)
File enforce strict encapsulation, but there is not export mechanism and partially extends access
Assembly fits all criteria except the export list. The export list is not explicitly specified, but is set using access modifiers.

$M(\text{C\#}) = 0$


\textbf{3) Interface Abstraction Capacity $I(L)$}

C\# provides interfaces, abstract classes and delegate for construction dedecited to behaviour abstraction
C\# provides class, struct, record and record struct for construction dedecited to behaviour abstraction

$I(\text{C\#}) = \frac{3}{4} = 0.75$


\textbf{4) Encapsulation Enforcement Coefficient $E(L)$}

С\# provide encapsulation bypass using: reflection and \texttt{unsafe} code, flag InternalsVisibleTo, Invoke mechanism:

$E(\text{C\#}) = 0$

\textbf{Final metric value:}
\begin{flalign*}
    & \qquad AE(\text{C\#}) = 0.2 \cdot V(\text{C\#}) + 0.25 \cdot M(\text{C\#}) + 0.25 \cdot I(\text{C\#}) + 0.3 \cdot E(\text{C\#}) &&\\
    & \qquad = 0.2 \cdot 1 + 0.25 \cdot 0 + 0.25 \cdot 0.75 + 0.3 \cdot 0 &&\\
    & \qquad = 0.3875 &&
\end{flalign*}


\subsubsection{Inheritance Semantics and Hierarchy Management}

\textbf{1) Hierarchy Simplicity Factor $H_s(L)$}

C\# supports only single class inheritance, but supports multiple interface inheritance:

\[
    H_s(\text{C\#}) = 1 - \frac{1 - 1}{M_{max} - 1} = 1
\]

\textbf{2) Linearization Determinism Score $L_d(L)$}

In C\#, class inheritance is strictly single (i.e., a class may inherit from exactly one base class), and interface implementation does not introduce ambiguity in method resolution because interfaces do not provide inherited implementation in the same sense (default interface methods are resolved with well-defined precedence rules). The specification defines a deterministic lookup process: member resolution proceeds along the class inheritance chain, and in the presence of interfaces, explicit rules govern conflict resolution.

$L_d(\text{C\#}) = 1$

\textbf{3) Constraint Protection Coefficient $C_p(L)$}

$p_1=1$: C\# provides explicit finalisation mechanisms via the \texttt{sealed} modifier for classes and methods, preventing further inheritance or overriding. 

$p_2=0$: while \texttt{sealed} prevents extension, C\# does not support fully closed hierarchies with enumerated subclasses (as in algebraic data types), thus lacking strict sealed hierarchy sets. 

$p_3=1$: C\# enforces explicit override semantics through the mandatory \texttt{override} keyword for redefining virtual members. 

$p_4=1$: methods are non-virtual by default and must be explicitly marked \texttt{virtual}, \texttt{abstract}, or \texttt{override} to participate in polymorphism. 

$p_5=1$: the language enforces a clear separation between abstract classes and interfaces, with distinct constraints on implementation and inheritance. 

$p_6=1$: the compiler enforces multiple inheritance restrictions (e.g., no multiple class inheritance, constraints on base class selection, and rules for interface implementation).

\[
    C_p(\text{C\#}) = \frac{1+0+1+1+1+1}{6} \approx 0.83
\]


\textbf{4) Fragile Base Mitigation Factor $F_b(L)$}

$s_1=1$: C\# requires explicit \texttt{override} annotations, guarding against inadvertent overriding.

$s_2=1$: strict access modifiers (e.g., \texttt{private/protected/internal}) control visibility and overridability of members. 

$s_3=0$: invocation of base class implementations via \texttt{base} is optional, so it does not enforce controlled superclass delegation. 

$s_4=1$: the language enforces variance (covariance and contravariance) rules in generics and method overriding. 

$s_5=1$: methods are non-virtual by default and requires explicit opt-in to any polymorphic behaviour. 

$s_6=1$: the compiler enforces strict override consistency, requiring signature matching and return type compatibility. 

\[
    F_b(\text{C\#}) = \frac{1+1+0+1+1+1}{6} = \frac{5}{6} \approx 0.83
\]

\textbf{Final metric value:}
\begin{flalign*}
    & \qquad ISRI(\text{C\#}) = 0.2 \cdot H_s(\text{C\#}) + 0.25 \cdot L_d(\text{C\#}) + 0.3 \cdot C_p(\text{C\#}) + 0.25 \cdot F_b(\text{C\#}) &&\\
    & \qquad = 0.2 \cdot 1 + 0.25 \cdot 1 + 0.3 \cdot \frac{5}{6} + 0.25 \cdot \frac{5}{6} &&\\
    & \qquad \approx 0.9083 &&
\end{flalign*}


\subsubsection{Composition and Alternatives to Inheritance}

\textbf{1) Class-Based Delegation Score $D_{cls}(L)$}

C\# does not provide native delegation constructs:
The delegate keyword exists in C\# \cite{ecma334}. It is a reference type that encapsulates a reference to a method with a specific signature. It is actively used for events, callbacks, and asynchronous patterns.
There is no built-in syntax for automatic forwarding in C\#. The developer must manually proxy the calls.

\[
    K_{del}(\text{C\#}) = 1, S_{fwd}(\text{C\#}) = 0 \therefore D_{cls}(\text{C\#}) = 0.5
\]

\textbf{2) Trait Expressiveness Score $T_{exp}(L)$}

When using interfaces with DIM, if a class implements two interfaces containing methods with the same signature and default implementation, the C\# compiler does not resolve the conflict automatically.
C\# does not allow an interface/trait to require the implementing class to have arbitrary methods, properties, or fields that are not included in the interface's contract.

$R_{conf}(\text{C\#}) = 0, O_{mod}(\text{C\#}) = 0 \therefore T_{exp}(\text{C\#}) = 0$

\textbf{3) Interface Composition Capacity $I_{comp}(L)$}

In C\#, a class can implement multiple interfaces, and there is no Strict upper limit. 
Starting with C\# 8.0, default interface methods have been introduced. Now the interface can contain: method implementations, static methods, private methods (can be reused inside the interface).

\[
    I_{comp}(\text{C\#}) = \frac{N_{impl}(\text{C\#}) + D_{def}(\text{C\#})}{2} = \frac{1 + 1}{2} = 1
\]

\textbf{Final metric value:}
\begin{flalign*}
    & \qquad CRFI(\text{C\#}) = 0.3 \cdot D_{cls}(\text{C\#}) + 0.35 \cdot T_{exp}(\text{C\#}) + 0.35 \cdot I_{comp}(\text{C\#}) &&\\
    & \qquad = 0.3 \cdot 0.5 + 0.35 \cdot 0 + 0.35 \cdot 1 &&\\
    & \qquad = 0.5 &&
\end{flalign*}


\subsubsection{Object Representation and Creation Model}

\textbf{1) Identity Semantics Score $S_{id}(L)$}

Type categories that can be declared in C\# \cite{ecma334}:
\begin{itemize}
    \item class - reference type
    \item struct - value type
    \item interface - reference type
    \item enum - value type
    \item delegate - reference type
\end{itemize}

\[
    S_{id}(\text{C\#}) = \frac{3}{5} = 0.6
\]

\textbf{2) Instantiation Paradigm Diversity $S_{cr}(L)$}

In the C\# language, four semantically different paradigms of object creation are distinguished:

\begin{itemize}
    \item $Constructor-based (new)$ defined by the specification as the main mechanism for creating instances of reference and meaningful types.
    \item $Literal initialisation$  allows the creation of objects without explicitly calling the constructor at the syntax level.
    \item \texttt{Copying instantiation} (\texttt{with} expressions for \texttt{record}) implementing non-destructive mutation and representing a separate functional creation model.
    \item $Stack memory (stackalloc)$ provides an alternative model for managing the lifetime of objects outside of heap allocation
\end{itemize}

Other mechanisms, such as object initializers, reflection, or \texttt{ICloneable}, do not form separate paradigms, since they are either syntactic extensions of \texttt{new}, or library level constructions rather than to linguistic semantics.

\[
    S_{cr}(\text{C\#}) = \frac{4}{4}
\]


\textbf{3) Meta-level Extensibility Score $S_{meta}(L)$}

The .NET platform specification provides advanced introspection through the \texttt{System.Reflection} namespace, which corresponds to level $I_{level}=1$, allowing the structure of types and objects to be queried at runtime. 
Additionally, the language allows limited modification of structures at runtime via \texttt{ExpandoObject}, dynamic typing (\texttt{dynamic}), and type generation through \texttt{System.Reflection.Emit}, which satisfies the criterion for level $I_{level}=2$. 
However, C\# does not support the meta-object protocol and does not allow the object creation process to be redefined (e.g., the semantics of the \texttt{new} operator), which precludes achieving level $3$. 

$I_{level}(\text{JS/Python/Ruby}) = 3$

\[
    S_{meta}(\text{C\#}) = \frac{2}{3} 
\]

\textbf{Final metric value:}
\begin{flalign*}
    & \qquad ORCI(\text{C\#}) = 0.35 \cdot S_{id}(\text{C\#}) + 0.3 \cdot S_{cr}(\text{C\#}) + 0.35 \cdot S_{meta}(\text{C\#}) &&\\
    & \qquad = 0.35 \cdot 0.6 + 0.3 \cdot 1 + 0.35 \cdot \frac{2}{3} &&\\
    & \qquad \approx 0.7433 &&
\end{flalign*}


\subsubsection{Polymorphism and Dispatch Mechanisms}

\textbf{1) Dispatch Scope Mechanism $M_{scope}(L)$}

C\# natively supports single dispatch via virtual method invocation. The method resolution depends on the runtime type of the receiver.

$I_{single}(\text{C\#})=1$

The language does not provide true multiple dispatch, because method overloading is resolved at compile time and does not consider the runtime types of all arguments.

$I_{multi}(\text{C\#})=0$ 

It supports a form of predicate-based dispatch through pattern-matching constructs (e.g., \texttt{switch}-expressions and \texttt{is}-patterns). This constructs enable branching based on both type and additional logical conditions.

$I_{pred}(\text{C\#})=1$

\[
    M_{scope}(\text{C\#}) = \frac{1 * 1 + 2 * 0 + 2 * 1}{5} = 0.6
\]

\textbf{2) Dispatch Safety Guarantee $S_{disp}(L)$}

C\# enforces compile-time verification of method dispatch correctness using static type checking. The method calls, the virtual dispatch resolution, and the overload selection must be type-safe and valid at compile time. The compiler is rejecting ambiguous or invalid dispatch scenarios.

$I_{static}(\text{C\#})=1$

In overriding virtual methods, it is not necessary to convert all possible derived types, and the pattern matching does not have to be exhaustive unless specific constructs (e.g., enums or sealed types) are used. The language does not provide complete definitions of dispatching.

$I_{exhaust}(\text{C\#})=0$

\[
    S_{disp}(\text{C\#}) = \frac{1 + 0}{2} = 0.5
\]

\textbf{3) Virtualisation Optimization Potential $O_{virt}(L)$}

First, C\# supports non-virtual methods by default and allows explicitly sealing overridden methods using the \texttt{sealed} modifier, satisfying the condition.

$K_{final}(\text{C\#})=1$

Second, the language provides sealed class hierarchies via the \texttt{sealed} keyword, preventing further inheritance and enabling more aggressive compile-time and JIT optimizations.

$K_{sealed}(\text{C\#})=1$. 

Third, C\# supports static dispatch through \texttt{static} methods, which are resolved at compile time, as well as recent extensions such as \texttt{static abstract} interface members that further strengthen compile-time polymorphism.

$K_{static}(\text{C\#})=1$ 

\[
O_{virt}(\text{C\#}) = 1
\]

\textbf{Final metric value:}
\begin{flalign*}
    & \qquad PDEI(\text{C\#}) = 0.35 \cdot M_{scope}(\text{C\#}) + 0.35 \cdot S_{disp}(\text{C\#}) + 0.3 \cdot O_{virt}(\text{C\#}) &&\\
    & \qquad = 0.35 \cdot 0.6 + 0.35 \cdot 0.5 + 0.3 \cdot 1 &&\\
    & \qquad = 0.685 &&
\end{flalign*}


\subsubsection{State Model, Mutability, and Concurrency Guarantees}

\textbf{1) Default Mutability Score $M_{def}(L)$}

All instances of reference and value types are mutable by default: fields can be modified after initialisation, and property accessors do not restrict modification without additional modifiers. Immutability is achieved exclusively through explicit language constructs, such as \texttt{readonly} for fields, \texttt{init} for properties, or the use of the \texttt{record} type, which confirms the absence of immutability by default at the language level. 

\[
    M_{def}(\text{C\#}) = 0
\]

\textbf{2) Static Reference Constraint Density $S_{dens}(L)$}

Reference-related constructs in C\# language include \cite{ecma334}: 
\begin{itemize}
    \item reference types (classes, interfaces, arrays, delegates)
    \item \texttt{ref}
    \item \texttt{out}
    \item \texttt{in}
    \item unsafe pointers (\texttt{*}, \texttt{\&})
    \item nullable reference types (\texttt{T?})
    \item \texttt{ref struct} (e.g., \texttt{Span<T>})
    \item \texttt{this}
    \item \texttt{base}
    \item \texttt{dynamic}
\end{itemize}

Among these, constructs imposing compile-time constraints: 
\begin{itemize}
    \item \texttt{ref} aliasing and definite assignment rules
    \item \texttt{out} mandatory assignment before return
    \item \texttt{in} (read-only enforcement)
    \item nullable reference types static - flow analysis for null safety
    \item \texttt{ref struct} stack-only lifetime and escape restrictions
    \item unsafe pointers - restricted to \texttt{unsafe} context
\end{itemize}

\[
    S_{dens}(\text{C\#}) = \frac{6}{10} = 0.6
\]

\textbf{3) Concurrency Primitive Safety $P_{conc}(L)$}

According to the C\# specification and .NET memory model, C\# does not provide compile-time guarantees of data race freedom. Instead, it relies on programmer discipline with constructs such as \texttt{lock}, \texttt{Monitor}, and \texttt{volatile}, without enforcing race freedom via the type system or runtime detection:

$S_{race}(\text{C\#}) = 0$

C\# provides structured concurrency through \texttt{Task}, \texttt{async}/\texttt{await}, and the Task Parallel Library, which encapsulate threads and support composable asynchronous workflows but do not enforce actor-style isolation or channel-based communication as core language primitives:

\[
    S_{model}(\text{C\#}) = 0.5
\]

\[
    P_{conc}(\text{C\#}) = 0.25
\]


\textbf{Final metric value:}
\begin{flalign*}
    & \qquad SCSI(\text{C\#}) = 0.35 \cdot M_{def}(\text{C\#}) + 0.25 \cdot S_{dens}(\text{C\#}) + 0.4 \cdot P_{conc}(\text{C\#}) &&\\
    & \qquad = 0.35 \cdot 0 + 0.25 \cdot 0.6 + 0.4 \cdot 0.25 &&\\
    & \qquad = 0.25 &&
\end{flalign*}


\subsubsection{Support for Formal Specification and Verifiability}

\textbf{1) Native Contract Constructs Score $C_{nat}(L)$}

C\# does not provide first-class syntactic constructs for designing-by-contract (preconditions, postconditions, invariants). Such functionality is provided via external mechanisms, which are not part of the core language syntax or semantics. 

\[
    n_{pre} = 0, n_{post} = 0, n_{inv} = 0 \Rightarrow C_{nat}(\text{C\#}) = 0
\]

\textbf{2) Verification Enforcement Level $V_{enf}(L)$}

C\# does not provide built-in static verification mechanisms capable of proving program correctness properties at compile time. Its type system is not dependent or refinement-based. However, the language and runtime support: extensive runtime checking, including type safety, array bounds checking, null reference checks, and optional contract-like validation via assertions or guard clauses, all enforced during execution by the CLR.

\[
    V_{enf}(\text{C\#}) = 0.5
\]

\textbf{Final metric value:}
\begin{flalign*}
    & \qquad FSVI(\text{C\#}) = 0.45 \cdot C_{nat}(\text{C\#}) + 0.55 \cdot V_{enf}(\text{C\#}) = 0.45 \cdot 0 + 0.55 \cdot 0.5 &&\\
    & \qquad = 0.275 &&
\end{flalign*}


\subsection{Metrics Evaluation for Java}

Java represents a statically typed, nominally structured object-oriented language with a formally specified type system and execution model defined in the Java Language Specification (JLS)~\cite{gosling2015java}. Its design emphasises type safety, backward compatibility, and controlled extensibility.

\subsubsection{Type System and Subtyping Model}

\textbf{1) Nominal-Structural Typing Score $S_{nom}(L)$}

The subtype relationships in Java are established through explicit declarations such as \texttt{extends} and \texttt{implements}. There is no structural equivalence subtyping in the language. All subtype relations must be declared in class or interface definitions, which sutisfy explicit subtype declarations. 

\[
    I_{struct}(\text{Java}) = 0, I_{nom}(\text{Java}) = 1 
    \Rightarrow S_{nom}(\text{Java}) = \frac{0 + 1}{2} = 0.5
\]


\textbf{2) Variance Formalization Score $S_{var}(L)$}

Java does not allow built-in variance (no explicit \texttt{in}/\texttt{out} annotations on type parameters), but supports use-site variance with wildcards (e.g., \texttt{? extends T}, \texttt{? super T}). Wildcards may be considered as limited and indirect form of variance. 

The JLS defines strict typing and subtyping rules, including wildcard capture semantics, to ensure type safety. The safety is achieved partly through restrictions: invariance of generic types and limited covariance support.

\[
    V_{exp}(\text{Java}) = 0.5, V_{sound}(\text{Java}) = 0.5  S_{var}(\text{Java}) = \frac{0.5 + 0.5}{2} = 0.5
\]


\textbf{3) Inheritance-Subtyping Separation Score $S_{rel}(L)$}

Java supports subtyping both through class inheritance and interface implementation. A class may become a subtype of one or more interfaces independently of extending another class. However, Java does not support arbitrary explicit subtype declarations between classes outside inheritance.

$I_{dec}(\text{Java}) = 0.5$

Java partially separates subtyping from implementation inheritance: subclassing (\texttt{extends}) creates subtype relations together with implementation inheritance, while interface implementation (\texttt{implements}) introduces subtype relations without requiring inheritance from another class.

$I_{ind}(\text{Java}) = 0.5$

The JLS distinguishes subclassing from subtyping at the level of types, including class types, interface types, and assignment compatibility. However, this distinction is not fully formalized as a separation between two independent relations in the specification.

$I_{form}(\text{Java}) = 0.5$

\[
    S_{rel}(\text{Java}) = \frac{0.5 + 0.5 + 0.5}{3} = 0.5
\]

\textbf{Final metric value:}
\begin{flalign*}
    & \qquad TSSI(\text{Java}) = 0.25 \cdot S_{nom}(\text{Java}) + 0.3 \cdot S_{var}(\text{Java}) + 0.45 \cdot S_{rel}(\text{Java}) &&\\
    & \qquad = 0.25 \cdot 0.5 + 0.3 \cdot 0.5 + 0.45 \cdot 0.5 &&\\
    & \qquad = 0.5 &&
\end{flalign*}


\subsubsection{Abstraction and Encapsulation Mechanisms}

\textbf{1) Visibility Granularity Index $V(L)$}

Java defines four access scopes: \texttt{private} (accessible only within the declaring class), default/package-private (accessible within the same package), \texttt{protected} (accessible within the package and by subclasses), and \texttt{public} (accessible from any context)\cite{gosling2015java}.

\[
    n_{levels}(\text{C\#})=7 \Rightarrow V(\text{Java}) = \frac{4}{7}
\]


\textbf{2) Module Isolation Strength $M(L)$}

The Java Platform Module System (JPMS) introduces explicit module boundaries with controlled exports. Java provides two primary modular mechanisms: packages, defined in JLS (\S7), which organize types without enforcing explicit export control, and modules, defined in the JPMS (\texttt{module-info.java}), which enforce strong encapsulation via explicit \texttt{exports} and \texttt{requires} directives. Therefore, module boundaries are independent from package organization, although explicit accessibility control is restricted to modules.

\[
    b_{strict}(\text{Java}) = 1, b_{total}(\text{Java}) = 2 \Rightarrow M(\text{Java}) = \frac{1}{2} = 0.5
\]


\textbf{3) Interface Abstraction Capacity $I(L)$}

Java provides two primary abstraction constructs: (i) \texttt{interface}, which defines pure behavioural contracts with method signatures (JLS \S9), and (ii) \texttt{abstract class}, which allows partial implementation alongside abstract methods (JLS \S8.1.1.1). The object-defining constructs in Java are classes records, enums and annotations.

\[
    n_{abstraction_constructs}(\text{Java}) = 2, n_{object}(\text{Java}) = 4 \Rightarrow I(\text{Java}) = \frac{2}{4} = 0.5
\]



\textbf{4) Encapsulation Enforcement Coefficient $E(L)$}

While Java enforces access control statically via modifiers such as \texttt{private}, \texttt{protected}, and \texttt{public} (JLS \S6.6), it also provides reflection capabilities through the core reflection API (\texttt{java.lang.reflect}), which allows runtime inspection and, critically, the ability to override access checks via mechanisms such as \texttt{setAccessible(true)} \cite{bloch2017effective}. Furthermore, low-level facilities (e.g., \texttt{Unsafe} and deep reflection) enable direct memory or field access outside standard type safety guarantees. 

\[
    E(\text{Java}) = 0
\]

\textbf{Final metric value:}
\begin{flalign*}
    & \qquad AE(\text{Java}) &&\\
    & \qquad = 0.2 \cdot V(\text{Java}) + 0.25 \cdot M(\text{Java}) + 0.25 \cdot I(\text{Java}) + 0.3 \cdot E(\text{Java}) &&\\
    & \qquad = 0.2 \cdot \frac{4}{7} + 0.25 \cdot 0.5 + 0.25 \cdot 0.5 + 0.3 \cdot 0 &&\\
    & \qquad = 0.3643 &&
\end{flalign*}


\subsubsection{Inheritance Semantics and Hierarchy Management}

\textbf{1) Hierarchy Simplicity Factor $H_s(L)$}

Java enforces single inheritance for classes, allowing each class to declare at most one direct superclass via the \texttt{extends} clause. Although Java supports multiple inheritance of type through interfaces, this does not affect $M(L)$, as the metric explicitly concerns superclass declarations.

\[
    H_s(\text{Java}) = 1
\]

\textbf{2) Linearization Determinism Score $L_d(L)$}

According to the Java Language Specification (JLS), method lookup is strictly defined through a well-ordered procedure: class inheritance follows a single inheritance model for classes, ensuring a unique linear superclass chain, while multiple inheritance of behaviour is permitted only via interfaces, whose default methods are resolved using explicitly defined rules (JLS \S8.4.8, \S9.4.1). In particular, Java enforces deterministic conflict resolution by prioritizing class methods over interface default methods, and requiring explicit disambiguation when multiple interfaces provide conflicting defaults.

\[
    L_d(\text{Java}) = 1
\]

\textbf{3) Constraint Protection Coefficient $C_p(L)$}

$p_1=1$: Java provides finalisation mechanism through the \texttt{final} keyword for classes and methods, preventing inheritance and overriding.

$p_2=1$: the language supports sealed classes and interfaces (standardised in Java~17), enabling explicitly restricted subclass hierarchies.

$p_3=0$: Java does not require explicit override declarations, since \texttt{@Override} annotation is optional.

$p_4=0$: methods are virtual by default (except when declared \texttt{final}, \texttt{private}, or \texttt{static}), so the language does not enforce non-virtual-by-default semantics.

$p_5=1$: Java maintains a clear syntactic and semantic distinction between abstract classes and interfaces, including separate inheritance and implementation rules.

$p_6=1$: the JLS defines compile-time inheritance restrictions, including the prohibition of multiple class inheritance, extensions of \texttt{final} classes, and absence of required abstract method implementations.

\[
    C_p(\text{Java}) = \frac{1+1+0+0+1+1}{6} \approx 0.67
\]

\textbf{4) Fragile Base Mitigation Factor $F_b(L)$}

$s_1=0$: Java does not require explicit override annotations, since \texttt{@Override} is optional.

$s_2=1$: the language enforces member visibility through access modifiers (\texttt{private}, \texttt{protected}, \texttt{public} and \texttt{package-private}), constraining inheritance and method overriding.

$s_3=1$: Java requires controlled superclass initialisation through explicit or implicit constructor invocation using \texttt{super}.

$s_4=1$: the JLS enforces type-safe overriding through return-type covariance and parameter invariance, preventing unsafe method redefinition.

$s_5=0$: methods are virtual by default (except when declaed \texttt{final}, \texttt{private} or \texttt{static}), so Java does not employ non-overridable-by-default semantics.

$s_6=1$: the compile performs strict compile-time checks for overriding consistency, including method signature matching, return type compatibility, and checked exception constraints. 

\[
    F_b(\text{Java}) = \frac{0+1+1+1+0+1}{6} \approx 0.67
\]

\textbf{Final metric value:}
\begin{flalign*}
    & \qquad ISRI(\text{Java}) &&\\
    & \qquad = 0.2 \cdot H_s(\text{Java}) + 0.25 \cdot L_d(\text{Java}) + 0.3 \cdot C_p(\text{Java}) + 0.25 \cdot F_b(\text{Java}) &&\\
    & \qquad = 0.2 \cdot \frac{4}{7} + 0.25 \cdot 1 + 0.3 \cdot \frac{4}{6} + 0.25 \cdot \frac{4}{6} &&\\
    & \qquad \approx 0.731 &&
\end{flalign*}


\subsubsection{Composition and Alternatives to Inheritance}

\textbf{1) Class-Based Delegation Score $D_{cls}(L)$}

Java does not provide built-in syntactic support for delegation mechanisms. Instead, delegation is typically implemented explicitly through object composition and manual method forwarding (often Adapter or Proxy pattern).

The Java Language Specification does not define automatic delegation or method-forwarding semantics. The language itself does not implicitly generate forwarding methods. So delegation behaviour is generally implemented manually or via external libraries and code-generation tools.

\[
    K_{del}(\text{Java}) = 0.5, S_{fwd}(\text{Java}) = 0 \Rightarrow D_{cls}(\text{Java}) = \frac{0.5 + 0}{2} = 0.25
\]


\textbf{2) Trait Expressiveness Score $T_{exp}(L)$}

In Java, conflicts arising from multiple inherited default methods using the same signatures cannot be resolved by aliasing or selective exclusion. The language causes a compile-time error if the derived class does not explicitly redefine the conflicting method and disambiguate it. 

Interfaces in Java cannot declare explicit implementation requirements beyond the requirements of method signatures. The language does not provide a mechanism for defining required fields, state dependencies and other structural constraints, or implementation-dependent behaviour in the interface itself.

\[
    R_{conf}(\text{Java}) = 0, O_{mod}(\text{Java}) = 0 \Rightarrow T_{exp}(\text{Java}) = 0
\]


\textbf{3) Interface Composition Capacity $I_{comp}(L)$}

Java allows a class to implement an arbitrary number of interfaces without any formal upper bound, supporting unrestricted interface composition. 

Since Java~8, interfaces may also define \textit{default methods} with concrete implementations, enabling interfaces to encapsulate reusable behaviour in addition to type abstraction.

\[
    N_{impl}(\text{Java}) = 1, D_{def}(\text{Java}) = 1 \Rightarrow I_{comp}(\text{Java}) = 1
\]


\textbf{Final metric value:}
\begin{flalign*}
    & \qquad CRFI(\text{Java}) = 0.3 \cdot D_{cls}(\text{Java}) + 0.35 \cdot T_{exp}(\text{Java}) + 0.35 \cdot I_{comp}(\text{Java}) &&\\
    & \qquad = 0.3 \cdot 0.25 + 0.35 \cdot 0 + 0.35 \cdot 1 &&\\
    & \qquad = 0.425 &&
\end{flalign*}


\subsubsection{Object Representation and Creation Model}

\textbf{1) Identity Semantics Score $S_{id}(L)$}

The main user-definable types are classes, interfaces, enums, records, and annotation types. They are all reference types, which means that variables hold references to objects rather than values, and identity is preserved via reference equality.

\[
    S_{id}(\text{Java}) = 1
\]

\textbf{2) Instantiation Paradigm Diversity $S_{cr}(L)$}

The primary instantiation mechanism is explicit construction via the \texttt{new} operator. Additionally, Java supports object creation via reflective instantiation (e.g., \texttt{Class::newInstance}, \texttt{Constructor::newInstance}). Furthermore, Java provides cloning through the \texttt{clone()} method. Other paradigms are not supported at the language level.

\[
    P_{supported}(\text{C\#}) = 4 \Rightarrow S_{cr}(\text{Java}) = \frac{3}{4}
\]

\textbf{3) Meta-level Extensibility Score $S_{meta}(L)$}

Java does not permit direct structural modification of loaded classes through the standard reflection API, although limited bytecode transformation is possible using instrumentation and specialised libraries.

\[
    I_{level}(\text{Java}) = 1, I_{level}(\text{JS/Python/Ruby}) = 3 \Rightarrow S_{meta}(\text{Java}) = \frac{1}{3}
\]


\textbf{Final metric value:}
\begin{flalign*}
    & \qquad ORCI(\text{Java}) = 0.35 \cdot S_{id}(\text{Java}) + 0.3 \cdot S_{cr}(\text{Java}) + 0.35 \cdot S_{meta}(\text{Java}) &&\\
    & \qquad = 0.35 \cdot 1 + 0.3 \cdot \frac{3}{4} + 0.35 \cdot \frac{1}{3} &&\\
    & \qquad \approx 0.6917 &&
\end{flalign*}


\subsubsection{Polymorphism and Dispatch Mechanisms}

\textbf{1) Dispatch Scope Mechanism $M_{scope}(L)$}

Java natively supports single dispatch, where method selection is based on the dynamic type of the receiver object.

$I_{single}(\text{Java}) = 1$. 

However, Java does not support multiple dispatch, because method overloading is resolved statically at compile time based on declared types, rather than dynamically across all arguments.

$I_{multi}(\text{Java}) = 0$

Similarly, predicate-based dispatch is not part of the language specification and must be emulated through explicit conditional logic.

$I_{pred}(\text{Java}) = 0$

\[
    M_{scope}(\text{Java}) = \frac{1}{5} = 0.2
\]

\textbf{2) Dispatch Safety Guarantee $S_{disp}(L)$}

Java uses static type checking to verify method invocations at compile time. Method calls are validated against declared reference types, overload resolution is performed statically during compilation, but overridden methods are dispatched dynamically at runtime.

$I_{static}(\text{Java}) = 1$

Java does not enforce exhaustiveness of dispatch definitions in general polymorphic scenarios (e.g., method overriding or dynamic dispatch across class hierarchies). Although newer features such as pattern matching in \texttt{switch} with sealed types introduce partial exhaustiveness checks, this is limited and not a general property of the dispatch system.

$I_{exhaust}(\text{Java}) = 0$

\[
    S_{disp}(\text{Java}) = \frac{1 + 0}{2} = 0.5
\]

\textbf{3) Virtualisation Optimization Potential $O_{virt}(L)$}

First, Java supports the \texttt{final} keyword for methods and classes, preventing overriding and enabling static binding.

$K_{final}(\text{Java}) = 1$

Second, since Java~17, the language includes \texttt{sealed} classes and interfaces, which restrict subclassing to a known finite set and allow more precise dispatch analysis.

$K_{sealed}(\text{Java}) = 1$

Third, Java provides static dispatch mechanisms through \texttt{static} methods and compile-time binding of non-virtual calls (e.g., \texttt{private}, \texttt{final} methods), which serve as explicit or implicit deVirtualisation hints.

$K_{static}(\text{Java}) = 1$

\[
    O_{virt}(\text{Java}) = 1
\]

\textbf{Final metric value:}
\begin{flalign*}
    & \qquad PDEI(\text{Java}) = 0.35 \cdot M_{scope}(\text{Java}) + 0.35 \cdot S_{disp}(\text{Java}) + 0.3 \cdot O_{virt}(\text{Java}) &&\\
    & \qquad = 0.35 \cdot 0.2 + 0.35 \cdot 0.5 + 0.3 \cdot 1 &&\\
    & \qquad = 0.545 &&
\end{flalign*}


\subsubsection{State Model, Mutability, and Concurrency Guarantees}

\textbf{1) Default Mutability Score $M_{def}(L)$}

In Java, all objects are mutable by default unless explicitly designed otherwise by the programmer, for example by declaring fields as \texttt{final} or creating immutable classes (e.g., \texttt{String}). The JLS specifies that instance fields are mutable unless restricted, and there is no language-level requirement or annotation enforcing immutability by default. Moreover, immutability in Java is a design pattern rather than a built-in default behaviour, requiring explicit developer effort. Therefore, Java does not satisfy the condition of immutability by default nor the condition where mutability requires explicit annotation. The mutability is inherent and unrestricted unless constrained.

\[
    M_{def}(\text{Java}) = 0
\]

\textbf{2) Static Reference Constraint Density $S_{dens}(L)$}

Java employs references for all non-primitive types, and the primary constructs contributing to $N_{ref}(L)$ include: object references (class types), array references, interface-typed references, generic type references, and the \texttt{null} reference.

Java employs references for all non-primitive types: object references (class types), array references, interface-typed references, generic type references, and the \texttt{null} reference.

$N_{ref}(\text{Java}) = 5$

Java provide constructs that impose compile-time constraints on reference usage: the \texttt{final} keyword, which restricts reassignment of references, generic type bounds (e.g., \texttt{<T extends ...>}), which constrain admissible reference types, and access modifiers (e.g., \texttt{private}, \texttt{protected}) that statically restrict reference visibility and usage contexts.

$N_{constr}(\text{Java}) = 3$

\[
    S_{dens}(\text{Java}) = 0.6
\]

\textbf{3) Concurrency Primitive Safety $P_{conc}(L)$}

First, Java does not provide compile-time guarantees of data race freedom. It relies on programmer discipline using constructs such as \texttt{synchronized}, \texttt{volatile}, and higher-level utilities from \texttt{java.util.concurrent}. Although the Java Memory Model formally defines happens-before relationships, it does not prevent data races statically nor enforce them via runtime checks \cite{buchholzjava}.

$S_{race}(\text{Java}) = 0$

Second, Java primarily exposes low-level concurrency primitives such as threads (\texttt{Thread}) and locks (\texttt{synchronized}, \texttt{Lock}), while also offering structured abstractions like \texttt{Future}, \texttt{ExecutorService}, and task-based concurrency. However, these do not constitute fully high-level models like Actors or CSP channels.

$S_{model}(\text{Java}) = 0.5$

\[
    P_{conc}(\text{Java}) = 0.25
\]


\textbf{Final metric value:}
\begin{flalign*}
    & \qquad SCSI(\text{Java}) = 0.35 \cdot M_{def}(\text{Java}) + 0.25 \cdot S_{dens}(\text{Java}) + 0.4 \cdot P_{conc}(\text{Java}) &&\\
    & \qquad = 0.35 \cdot 0 + 0.25 \cdot 0.6 + 0.4 \cdot 0.25 = 0.25 &&
\end{flalign*}


\subsubsection{Support for Formal Specification and Verifiability}

\textbf{1) Native Contract Constructs Score $C_{nat}(L)$}

Java does not include native language-level constructs for specifying preconditions, postconditions, or class invariants as part of its syntax or type system. While developers can express such constraints using assertions (\texttt{assert}) or external frameworks, these are not integrated as formal contract constructs with guaranteed enforcement semantics in the core language specification.

\[
    C_{nat}(\text{Java}) = 0
\]

\textbf{2) Verification Enforcement Level $V_{enf}(L)$}

Java does not support static verification of program correctness in the sense of compile-time proofs for contracts or general behavioural properties. The type system ensures type safety but does not verify semantic correctness such as preconditions or invariants. Java mostly relies on runtime assertions and exception-based validation checking mechanisms to detect violations of expected program behaviour during execution. These mechanisms are explicitly defined in the language and standard runtime but operate only at runtime rather than compile time.

\[
    V_{enf}(\text{Java}) = 0.5
\]

\textbf{Final metric value:}
\begin{flalign*}
    & \qquad FSVI(\text{Java}) = 0.45 \cdot C_{nat}(\text{Java}) + 0.55 \cdot V_{enf}(\text{Java}) = 0.45 \cdot 0 + 0.55 \cdot 0.5 &&\\
    & \qquad = 0.275 &&
\end{flalign*}


\subsection{Metrics Evaluation for Smalltalk}

Smalltalk represents one of the earliest and most influential realizations of a pure object-oriented programming paradigm, where everything is an object and computation is performed exclusively via message passing \cite{goldberg1983smalltalk}. The following chapter evaluates Smalltalk against the proposed metrics, grounding each numerical assignment in the formal and de facto language specification as documented in canonical literature and subsequent analyses.

\subsubsection{Type System and Subtyping Model}

\textbf{1) Nominal-Structural Typing Score $S_{nom}(L)$}

Smalltalk is a purely dynamically typed language\cite{blackpharo} whose semantics are defined by message passing rather than static typing judgments. It does not provide formal structural subtyping rules.

$I_{struct}(\text{Smalltalk}) = 0$

Smalltalk subtype relationships are not enforced through a static type system with explicit nominal typing judgments, but instead emerge dynamically at runtime. It is lack of formal nominal typing as defined by static type theory.

$I_{nom}(\text{Smalltalk}) = 0$

\[
    S_{nom}(\text{Smalltalk}) = \frac{0 + 0}{2} = 0
\]

\textbf{2) Variance Formalization Score $S_{var}(L)$}

Smalltalk does not include generics or parametric type constructors in its core language definition. All polymorphism is achieved dynamically via message passing without static type parameters, hence no variance annotations (covariance, contravariance, or invariance) can be expressed.

$V_{exp}(\text{Smalltalk}) = 0$

Additionally, since variance is not part of the language semantics, there are no formally defined typing rules or proofs of type soundness related to variance in the specification.

$V_{sound}(\text{Smalltalk}) = 0$

\[
    V_{exp}(\text{Smalltalk}) = 0, V_{sound}(\text{Smalltalk}) = 0 \Rightarrow S_{var}(\text{Smalltalk}) = 0
\]

\textbf{3) Inheritance-Subtyping Separation Score $S_{rel}(L)$}

First, Smalltalk does not provide any explicit notion of subtyping separate from class inheritance. The type compatibility is entirely governed by the class hierarchy, and there are no interfaces, traits, or structural typing mechanisms that would allow subtyping without inheritance.

$I_{dec}(\text{Smalltalk})=0$

Second, the subtype relation is not defined independently but is implicitly identical to the inheritance relation: an object is considered substitutable only if it belongs to a subclass, reflecting purely nominal and inheritance-based typing.

$I_{ind}(\text{Smalltalk})=0$

Third, the specification does not formally distinguish between inheritance and subtyping as separate relations. The behavioural compatibility is informally tied to method availability within the class hierarchy, without a formal type system that separates these concepts.

$I_{form}(\text{Smalltalk})=0$

\[
    S_{rel}(\text{Smalltalk}) = \frac{0+0+0}{3} = 0
\]

\textbf{Final metric value:}
\begin{flalign*}
    & \qquad TSSI(\text{Smalltalk}) &&\\
    & \qquad = 0.3 \cdot S_{nom}(\text{Smalltalk}) + 0.3 \cdot S_{var}(\text{Smalltalk}) + 0.4 \cdot S_{rel}(\text{Smalltalk}) &&\\
    & \qquad = 0.3 \cdot 0 + 0.3 \cdot 0 + 0.4 \cdot 0 = 0 &&
\end{flalign*}

\begin{flalign*}
    & \qquad TSSI(\text{Smalltalk}) &&\\
    & \qquad = 0.25 \cdot S_{nom}(\text{Smalltalk}) + 0.3 \cdot S_{var}(\text{Smalltalk}) + 0.45 \cdot S_{rel}(\text{Smalltalk}) &&\\
    & \qquad = 0.25 \cdot 0 + 0.3 \cdot 0 + 0.45 \cdot 0 = 0 &&
\end{flalign*}


\subsubsection{Abstraction and Encapsulation Mechanisms}

\textbf{1) Visibility Granularity Index $V(L)$}

Smalltalk does not provide formal access modifiers such as \textit{private}, \textit{protected}, or module-level visibility. All methods are publicly accessible via message sending, and encapsulation is achieved by convention (e.g., method categories) rather than enforced by the language semantics. 

$n_{levels}(\text{Smalltalk}) = 1$

\[
    n_{levels}(\text{C\#})=7 \Rightarrow V(\text{Smalltalk}) = \frac{1}{7}
\]


\textbf{2) Module Isolation Strength $M(L)$}

Smalltalk organizes code through classes and method categories but lacks a strict module system with enforced visibility or explicit export control. Classes and methods mostly are globally accessible through message passing.

\[
    M(\text{Smalltalk}) = 0
\]


\textbf{3) Interface Abstraction Capacity $I(L)$}

Smalltalk defines objects exclusively through classes.

$n_{object}(\text{Smalltalk}) = 1$

The language lacks interfaces, enforced abstract classes, and formal protocols. Behavioural abstraction is achieved through duck typing and message passing.

$n_{abstraction\_constructs}(\text{Smalltalk}) = 0$

\[
    
    I(\text{Smalltalk}) = \frac{0}{1} = 0
\]


\textbf{4) Encapsulation Enforcement Coefficient $E(L)$}

Smalltalk supports full reflection and allows runtime access and modification of objects, classes, and object internals without enforced encapsulation restrictions.

\[
    E(\text{Smalltalk}) = 0
\]

\textbf{Final metric value:}
\begin{flalign*}
    & \qquad AE(\text{Smalltalk}) &&\\
    & \qquad = 0.2 \cdot V(\text{Smalltalk}) + 0.25 \cdot M(\text{Smalltalk}) + 0.25 \cdot I(\text{Smalltalk}) &&\\
    & \qquad \qquad + 0.3 \cdot E(\text{Smalltalk}) &&\\
    & \qquad = 0.2 \cdot \frac{1}{7} + 0.25 \cdot 0 + 0.25 \cdot 0 + 0.3 \cdot 0 &&\\
    & \qquad \approx 0.0286 &&
\end{flalign*}


\subsubsection{Inheritance Semantics and Hierarchy Management}

\textbf{1) Hierarchy Simplicity Factor $H_s(L)$}

Smalltalk enforces a \emph{single inheritance} paradigm, meaning that each class has exactly one direct superclass (except for the root class, typically \texttt{Object}).

\[
    H_s(\text{Smalltalk}) = 1
\]

\textbf{2) Linearization Determinism Score $L_d(L)$}

Smalltalk employs a strictly single inheritance model, where each class has exactly one superclass, forming a linear inheritance chain. Method resolution is therefore performed via a deterministic upward traversal of this chain \cite{goldberg1983smalltalk}, starting from the receiver’s class and proceeding through its superclasses until a matching method is found.

\[
    L_d(\text{Smalltalk})=1
\]

\textbf{3) Constraint Protection Coefficient $C_p(L)$}

$p_1=0$: Smalltalk does not support final classes or methods.

$p_2=0$: there is no notion of sealed or closed hierarchies, and any class can be extended freely.

$p_3=0$: as method overriding does not require explicit annotation (method replacement is implicit via dynamic dispatch).

$p_4=0$: all methods are virtual by default.

$p_5=0$: the language lacks formal abstract classes and interfaces.

$p_6=0$: Smalltalk has no compile-time inheritance restrictions 

\[
    C_p(\text{Smalltalk})=\frac{0+0+0+0+0+0}{6}=0
\]

\textbf{4) Fragile Base Mitigation Factor $F_b(L)$}

$s_1=0$: Overriding does not require explicit annotation and occurs implicitly.

$s_2=0$: Method visibility is not strictly enforced.

$s_3=0$: Superclass method invocation is optional.

$s_4=0$: Smalltalk does not define variance constraints.

$s_5=0$: All methods are overridable by default.

$s_6=0$: Smalltalk has no compile-time override checks. 

\[
    F_b(\text{Smalltalk})=\frac{0+0+0+0+0+0}{6}=0
\]

\textbf{Final metric value:}
\begin{flalign*}
    & \qquad ISRI(\text{Smalltalk}) &&\\
    & \qquad = 0.2 \cdot H_s(\text{Smalltalk}) + 0.25 \cdot L_d(\text{Smalltalk}) + 0.3 \cdot C_p(\text{Smalltalk}) &&\\
    & \qquad \qquad + 0.25 \cdot F_b(\text{Smalltalk}) &&\\
    & \qquad = 0.2 \cdot 1 + 0.25 \cdot 1 + 0.3 \cdot 0 + 0.25 \cdot 0 &&\\
    & \qquad = 0.45 &&
\end{flalign*}


\subsubsection{Composition and Alternatives to Inheritance}

\textbf{1) Class-Based Delegation Score $D_{cls}(L)$}

$K_{del}(\text{Smalltalk})=0.5$: Smalltalk does not provide built-in keywords for delegation, but supports delegation through composition and message forwarding via \texttt{doesNotUnderstand:}.

$S_{fwd}(\text{Smalltalk})=1$: Smalltalk supports automatic delegation through \texttt{doesNotUnderstand:}.

\[
    D_{cls}(\text{Smalltalk})=0.75
\]

\textbf{2) Trait Expressiveness Score $T_{exp}(L)$}

Smalltalk traits provide explicit mechanisms for resolving method name conflicts through both exclusion and aliasing operators, which are formally defined in the trait composition semantics.
The traits can declare required methods that must be implemented by the consuming class, enabling explicit specification of behavioural dependencies.

\[
    R_{conf}(\text{Smalltalk})=1, O_{mod}(\text{Smalltalk})=1 \Rightarrow T_{exp}(\text{Smalltalk})=0
\]

\textbf{3) Interface Composition Capacity $I_{comp}(L)$}

Smalltalk does not define interfaces as a separate language feature. It relies on implicit typing (duck typing) and behavioural conformance, where any object responding to a set of messages satisfies the expected protocol. Consequently, there are no formal interface declarations or constructs analogous to interfaces, yielding Since interfaces do not exist as entities, there is no notion of default method implementations within them, thus 

\[
     N_{impl}(\text{Smalltalk})=0, D_{def}(\text{Smalltalk})=0 \Rightarrow I_{comp}(\text{Smalltalk})=0
\]

\textbf{Final metric value:}
\begin{flalign*}
    & \qquad CRFI(\text{Smalltalk}) &&\\
    & \qquad = 0.3 \cdot D_{cls}(\text{Smalltalk}) + 0.35 \cdot T_{exp}(\text{Smalltalk}) + 0.35 \cdot I_{comp}(\text{Smalltalk}) &&\\
    & \qquad = 0.3 \cdot 0.75 + 0.35 \cdot 0 + 0.35 \cdot 0 &&\\
    & \qquad = 0.225 &&
\end{flalign*}


\subsubsection{Object Representation and Creation Model}

\textbf{1) Identity Semantics Score $S_{id}(L)$}

In Smalltalk, all user-definable types are classes, and every instance is an object accessed via references, with no distinction between value types and reference types at the language level.

\[
    S_{id}(\text{Smalltalk})=1
\]

\textbf{2) Instantiation Paradigm Diversity $S_{cr}(L)$}

Smalltalk natively supports one primary instantiation paradigm: class-based object creation via \texttt{new}

\[
    P_{supported}(\text{C\#}) = 4 \Rightarrow S_{cr}(\text{Smalltalk})=\frac{1}{4}
\]

\textbf{3) Meta-level Extensibility Score $S_{meta}(L)$}

Smalltalk is a fully reflective environment where classes are objects (instances of metaclasses), and their structure and behaviour can be modified at runtime. Object creation can be customized by overriding class-side methods such as \texttt{new}, \texttt{basicNew}, and initialisation protocols.
\begin{flalign*}
    & \qquad I_{level}(\text{Smalltalk})=3, I_{level}(\text{JS/Python/Ruby}) = 3 &&\\
    & \qquad \Rightarrow S_{meta}(\text{Smalltalk})= \frac{1}{3} &&
\end{flalign*}


\textbf{Final metric value:}
\begin{flalign*}
    & \qquad ORCI(\text{Smalltalk}) = 0.35 \cdot S_{id}(\text{Smalltalk}) + 0.3 \cdot S_{cr}(\text{Smalltalk}) &&\\
    & \qquad \qquad + 0.35 \cdot S_{meta}(\text{Smalltalk}) &&\\
    & \qquad = 0.35 \cdot 1 + 0.3 \cdot \frac{1}{4} + 0.35 \cdot \frac{1}{3} &&\\
    & \qquad \approx 0.5417 &&
\end{flalign*}


\subsubsection{Polymorphism and Dispatch Mechanisms}

\textbf{1) Dispatch Scope Mechanism $M_{scope}(L)$}

Smalltalk natively supports single dispatch, where method selection depends solely on the dynamic type of the message receiver.

$I_{single}(\text{Smalltalk})=1$

However, it does not provide native multiple dispatch, as method resolution does not consider the dynamic types of arguments beyond the receiver.

$I_{multi}(\text{Smalltalk})=0$

In Smalltalk, conditional behaviour is implemented inside method bodies rather than through the dispatch mechanism. They are not part of language itself.

$I_{pred}(\text{Smalltalk})=0$ 

\[
    M_{scope}(\text{Smalltalk})=\frac{1}{5}=0.2
\]

\textbf{2) Dispatch Safety Guarantee $S_{disp}(L)$}

First, Smalltalk employs purely dynamic message passing with no static type checking. The correctness of a message send (i.e., whether the receiver understands the selector) is verified only at runtime, potentially resulting in a \texttt{doesNotUnderstand:} exception.

$I_{static}(\text{Smalltalk})=0$

Second, the language does not require exhaustive definitions of dispatch cases-method lookup proceeds dynamically through the class hierarchy, and missing methods are handled at runtime rather than enforced statically.

$I_{exhaust}(\text{Smalltalk})=0$

\[
    S_{disp}(\text{Smalltalk})=\frac{0+0}{2}=0
\]

\textbf{3) Virtualisation Optimization Potential $O_{virt}(L)$}

Smalltalk does not support \texttt{final} or non-virtual methods, as all methods are dynamically dispatched and can be overridden.

$K_{final}(\text{Smalltalk})=0$

Similarly, it does not provide sealed class hierarchies or mechanisms to restrict subclassing.

$K_{sealed}(\text{Smalltalk})=0$

Furthermore, there are no static dispatch hints or constructs such as \texttt{static} methods or annotations for deVirtualisation. All message sends are uniformly dynamic.

$K_{static}(\text{Smalltalk})=0$

\[
    O_{virt}(\text{Smalltalk})=0
\]

\textbf{Final metric value:}
\begin{flalign*}
    & \qquad PDEI(\text{Smalltalk}) = 0.35 \cdot M_{scope}(\text{Smalltalk}) + 0.35 \cdot S_{disp}(\text{Smalltalk}) &&\\
    & \qquad \qquad + 0.3 \cdot O_{virt}(\text{Smalltalk}) &&\\
    & \qquad = 0.35 \cdot 0.2 + 0.35 \cdot 0 + 0.3 \cdot 0 &&\\
    & \qquad = 0.07 &&
\end{flalign*}


\subsubsection{State Model, Mutability, and Concurrency Guarantees}

\textbf{1) Default Mutability Score $M_{def}(L)$}

In Smalltalk, all entities are objects and their state is inherently mutable: instance variables can be freely modified through message passing without requiring any special annotations or qualifiers. There is no built-in notion of immutability enforced by default, nor a requirement to explicitly mark objects as mutable, mutability is the standard behaviour.

\[
    M_{def}(\text{Smalltalk})=0
\]

\textbf{2) Static Reference Constraint Density $S_{dens}(L)$}

In Smalltalk, all variables (including instance, temporary, and class variables) uniformly denote references to objects, and message passing is the sole mechanism for interaction. The distinct reference-related constructs include: variable bindings, assignment operation, and message sends that operate on object references.

$N_{ref}(\text{Smalltalk}) = 3$

However, Smalltalk is a dynamically typed language with no static type system and no compile-time enforcement of aliasing, mutability, or reference usage constraints, all such checks occur at runtime. Consequently, none of the reference-related constructs impose static constraints.

$N_{constr}(\text{Smalltalk}) = 0$

\[
    S_{dens}(\text{Smalltalk})=0
\]

\textbf{3) Concurrency Primitive Safety $P_{conc}(L)$}

First, Smalltalk provides lightweight processes (green threads) and synchronisation primitives (e.g., semaphores), but lacks a static type system or runtime mechanisms that enforce data race freedom, correctness depends entirely on programmer discipline when coordinating shared mutable state.

$S_{race}(\text{Smalltalk}) = 0$

Second, Smalltalk offers low-level concurrency constructs such as processes and semaphores rather than structured concurrency or high-level models like Actors or CSP channels in its standard specification.

$S_{model}(\text{Smalltalk}) = 0$

\[
    P_{conc}(\text{Smalltalk})=0
\]

\textbf{Final metric value:}
\begin{flalign*}
    & \qquad SCSI(\text{Smalltalk}) &&\\
    & \qquad = 0.35 \cdot M_{def}(\text{Smalltalk}) + 0.25 \cdot S_{dens}(\text{Smalltalk}) &&\\
    & \qquad \qquad + 0.4 \cdot P_{conc}(\text{Smalltalk}) &&\\
    & \qquad = 0.35 \cdot 0 + 0.25 \cdot 0 + 0.4 \cdot 0 = 0 &&
\end{flalign*}


\subsubsection{8) Support for Formal Specification and Verifiability}

\textbf{1) Native Contract Constructs Score $C_{nat}(L)$}

Smalltalk does not natively support preconditions, postconditions, or invariants as part of its syntax or semantics, such checks must be implemented manually using message sends (e.g., assertions) or via external frameworks, and are not enforced by the language itself.

\[
    C_{nat}(\text{Smalltalk})=0
\]

\textbf{2) Verification Enforcement Level $V_{enf}(L)$}

Smalltalk is a dynamically typed language without built-in support for static verification or compile-time correctness proofs. Runtime checks can be implemented manually.

\[
    V_{enf}(\text{Smalltalk})=0
\]

\textbf{Final metric value:}
\begin{flalign*}
    & \qquad FSVI(\text{Smalltalk}) = 0.45 \cdot C_{nat}(\text{Smalltalk}) + 0.55 \cdot V_{enf}(\text{Smalltalk}) &&\\
    & \qquad = 0.45 \cdot 0 + 0.55 \cdot 0 = 0 &&
\end{flalign*}

